\section{Additional Results}

\noindent{\textbf{Video 1.}} To provide further intuition, we additionally provide videos on the project page that visualize the spatial distribution of allocated Gaussians under varying Gaussian budgets for the samples shown in~\cref{fig:teaser} and \cref{fig:densification_score_analysis} of the main paper. In \cref{fig:densification_score_analysis} of the main paper, for clearer visualization, only the locations allocated at the finest level ($l=3$) are highlighted in red. In the supplementary videos, by contrast, all allocated locations are shown in red. These videos illustrate that, as the Gaussian budget increases, Gaussians are adaptively allocated to fine-detail regions, and redundant allocations are also minimized in regions with overlap among the context images.

\noindent{\textbf{Video 2.}} Video 2 on the project page provides comparisons with AnySplat~\cite{jiang2025anysplat} under different Gaussian budgets. Both models use weights trained on the multi-view RE10K setting, and the videos are generated by taking two input views from the RE10K and ACID datasets. For AnySplat, we performed inference while progressively doubling the voxel size from its default value. We then used the number of Gaussians produced by AnySplat as the Gaussian budget for \OURMODEL. The first row shows the rendered RGB quality, and the second row presents the corresponding depth maps. In the third row, allocated Gaussian locations are highlighted in red. While AnySplat distributes Gaussians uniformly over the scene, \OURMODEL\ allocates them in a spatially adaptive manner. Ultimately, \OURMODEL\ maintains high fidelity even with fewer Gaussians.


\noindent{\textbf{\cref{fig:supple_qual_re10k}.}} We also provide additional qualitative comparisons for the multi-view experiments in \cref{fig:supple_qual_re10k}.